\section{A larger explicit family: sums of binary
forms}\label{a-larger-explicit-family-sums-of-binary-forms}

\textbf{Scope.} This is a further application of the finite rational
certificate theorem CC-XC05--07 and the geometric closure theorem
CC-XC08, with the same explicitly stated abelian-sector inputs. The
source geometry is Shioda--Katsura (1979, Remark 1.10, p.~102), which
extends their construction from Fermat equations to sums of homogeneous
forms in disjoint variables. A direct verification of the part of that
construction needed here is included. No claim of novelty or of the Tate
conjecture for arbitrary hypersurfaces is made.

\subsection{CC-XC10 --- Separated-variable joins and guaranteed
termination}\label{cc-xc10-separated-variable-joins-and-guaranteed-termination}

Let \(k=\mathbf F_q\) have characteristic \(p\), and let \(m\ge2\) with
\(p\nmid m\).

\textbf{Join criterion.} Let \(r,s\ge1\), and let

\[
 f(x_0,\ldots,x_r),\qquad g(y_0,\ldots,y_s)
\]

be homogeneous polynomials of degree \(m\). Suppose the projective
hypersurfaces

\[
 F=\{f=0\}\subset\mathbf P^r,\qquad
 G=\{g=0\}\subset\mathbf P^s
\]

are smooth. Define

\[
 X_f=\{f(x)+u^m=0\},\quad
 X_g=\{g(y)+v^m=0\},\quad
 Z_{f,g}=\{f(x)+g(y)=0\}.
 \tag{CC-XC.S24}
\]

These are smooth projective hypersurfaces of dimensions \(r,s,r+s\). If
\(F,G,X_f,X_g\) are certified in the sense of CC-XC08, then \(Z_{f,g}\)
is certified.

\textbf{Explicit family.} For every \(t\ge1\), every collection of
nonzero squarefree homogeneous binary forms \(f_1,\ldots,f_t\) of degree
\(m\) over \(k\), and every \(a\in k^\times\), both hypersurfaces

\[
 \begin{aligned}
 B_t(f_1,\ldots,f_t)
   &:\ \sum_{j=1}^t f_j(x_{2j-2},x_{2j-1})=0
       &&\subset\mathbf P^{2t-1},\\
 A_t(f_1,\ldots,f_t;a)
   &:\ \sum_{j=1}^t f_j(x_{2j-2},x_{2j-1})+a z^m=0
       &&\subset\mathbf P^{2t}
 \end{aligned}
 \tag{CC-XC.S25}
\]

are smooth and certified. Here squarefree means no repeated linear
factor after extending to \(\bar k\); equivalently the corresponding
degree-\(m\) subscheme of \(\mathbf P^1\) is finite étale. In particular
the finite certificate search of CC-XC07 is guaranteed to halt on every
input in (CC-XC.S25). These include hypersurfaces whose binary blocks
have arbitrary distinct branch points, rather than only diagonal
monomials.

\textbf{Complete proof of the join criterion.}

First, the affine gradients of \(f\) and \(g\) vanish only at the
origin. Indeed, if all first partials of \(f\) vanish at a nonzero
vector \(x\), Euler's identity gives

\[
 m f(x)=\sum_i x_i\frac{\partial f}{\partial x_i}(x)=0.
\]

Since \(p\nmid m\), this would be a singular point of the projective
hypersurface \(F\). The same reasoning applies to \(g\). It follows at
once that every hypersurface in (CC-XC.S24) is smooth: vanishing of its
partial derivatives would force all its projective coordinates to
vanish. Their positive dimensions imply geometric integrality, by the
component-intersection argument in CC-XC09.

Choose a finite extension containing \(\epsilon\) with
\(\epsilon^m=-1\). It is enough to prove certification there, by
CC-XC08(4). For clarity use \(k\) for this enlarged field during the
construction. The map

\[
 \phi:X_f\times X_g\dashrightarrow Z_{f,g},\qquad
 ((x:u),(y:v))\longmapsto(x_0v:\cdots:x_rv:
                              \epsilon u y_0:\cdots:\epsilon u y_s)
 \tag{CC-XC.S26}
\]

lands in the stated target because

\[
 f(xv)+g(\epsilon uy)
 =v^m(-u^m)+(-1)u^m(-v^m)=0.
\]

At a point of \(X_f\), not all the \(x_i\) vanish, since that would also
force \(u=0\). Similarly, some \(y_j\) is nonzero on \(X_g\). On the
product chart \(x_i y_j\ne0\), the ideal generated by the coordinates in
(CC-XC.S26) is therefore exactly the ideal generated by \(u/x_i\) and
\(v/y_j\), up to the line-bundle trivialization. Thus the base scheme is

\[
 C=F\times G=\{u=v=0\}\subset X_f\times X_g.
\]

It is smooth of codimension two. Blowing up its ideal makes this
generating linear system invertible, and hence resolves the rational map
to projective space. The Fermat-type equation already verified on the
dense domain of \(\phi\) then holds on the blowup. We obtain a proper
morphism

\[
 \widetilde\phi:\operatorname{Bl}_{F\times G}(X_f\times X_g)
                     \longrightarrow Z_{f,g}.
 \tag{CC-XC.S27}
\]

This is exactly the portion of Shioda--Katsura, Remark 1.10, p.~102,
required here; the preceding ideal calculation verifies the resolution
directly.

We also verify its degree. The exceptional divisor maps into
\(\{f(z_x)=g(z_y)=0\}\): over \((x,y)\in F\times G\), the leading
coordinates of the linear system have their two blocks proportional to
\(x\) and \(y\). Points of the original domain with \(u=0\) or \(v=0\)
also map into this closed set. Take a geometric point represented as
\((z_x:z_y)\) in the dense open where \(f(z_x)\ne0\). Then
\(g(z_y)=-f(z_x)\), so both coordinate blocks are nonzero, and its
entire inverse image lies in the original domain with \(uv\ne0\). Each
root

\[
 u^m=-f(z_x)
\]

gives a preimage

\[
 ((z_x:u),(z_y/(\epsilon u):1)).
 \tag{CC-XC.S28}
\]

The first point lies on \(X_f\), and

\[
 g(z_y/(\epsilon u))
 =\frac{-f(z_x)}{(-1)(-f(z_x))}=-1,
\]

so the second lies on \(X_g\). Conversely, in every preimage of this
open point, both \(u\) and \(v\) are nonzero. Normalize \(v=1\), and
then use the scalar of the first projective factor to make its
\(x\)-block equal to \(z_x\). The equations force exactly the
presentation (CC-XC.S28). More precisely, on a target chart where one
coordinate in \(z_x\) is normalized to \(1\), these formulas and the
normalizations are regular inverse maps identifying the inverse image
with the cover defined by \(u^m+f(z_x)=0\). This is finite étale of
degree \(m\), because \(p\nmid m\) and \(f(z_x)\) is invertible on the
chosen open. Thus the morphism is dominant and generically finite of
degree \(m\).

Since \(X_f,X_g,F,G\) are certified, so are their two products.
CC-XC08(3) certifies the blowup in (CC-XC.S27), and CC-XC08(2) certifies
its target by the degree-\(m\) pullback/trace splitting. Finally
CC-XC08(4) descends this conclusion to the original field. This proves
the join criterion.

\textbf{Complete proof for the binary-form family.}

For each squarefree binary \(f_j\), its gradient has no nonzero common
zero by the same Euler argument: a nonzero common zero would be a
repeated point of its degree-\(m\) divisor on \(\mathbf P^1\). As the
blocks involve disjoint variables, the gradient of \(\sum f_j\) vanishes
only when every block is zero. Adding \(a z^m\) preserves this property,
since \(ma\ne0\). Thus both hypersurfaces in (CC-XC.S25) are smooth.

We prove simultaneously, by induction on \(t\), the assertions for all
choices of the binary forms and all \(a\ne0\). For \(t=1\), \(B_1\) is
finite étale and \(A_1\) is a smooth projective curve, so CC-XC08(1)
applies.

Suppose \(t\ge2\) and all assertions with \(t-1\) blocks have been
proved. To prove the assertion for \(B_t\), put

\[
 f=\sum_{j<t}f_j,\qquad g=f_t.
\]

The four varieties in the join criterion are respectively

\[
 F=B_{t-1},\qquad G=B_1(f_t),\qquad
 X_f=A_{t-1}(f_1,\ldots,f_{t-1};1),\qquad
 X_g=A_1(f_t;1).
 \tag{CC-XC.S29}
\]

All are certified by induction or the base case. Their dimensions fit
the join criterion with \(r=2t-3\ge1\) and \(s=1\), and its target is
\(B_t\). This proves the \(B_t\) assertion for every collection of \(t\)
squarefree binary forms.

Now prove the \(A_t\) assertion. Put

\[
 f=\sum_{j<t}f_j+a z^m,\qquad g=f_t.
\]

Here the four input varieties are

\[
 \begin{aligned}
 F&=A_{t-1}(f_1,\ldots,f_{t-1};a),&
 G&=B_1(f_t),\\
 X_f&=B_t(f_1,\ldots,f_{t-1},h),&
 X_g&=A_1(f_t;1),
 \end{aligned}
 \qquad h(z,u)=a z^m+u^m.
 \tag{CC-XC.S30}
\]

The binary form \(h\) is squarefree because \(p\nmid m\) and \(a\ne0\).
Thus \(X_f\) is certified by the \(B_t\) assertion just proved, not by
an unproved \(A_t\) induction hypothesis. The other three inputs are
certified by induction or the base case. The dimensions are
\(r=2t-2\ge2\) and \(s=1\), and the join target is \(A_t\). This
completes the simultaneous induction without a circular step.

Certification gives all the rational, prime-independent splitting and
numerical-category consequences in CC-XC06. The positive search of
CC-XC07 therefore terminates on every member of the stated family. This
is a specified geometric class with a finite inductive construction; it
does not give a termination theorem for a general homogeneous
polynomial. \(\square\)

\subsection{Exact primary reference}\label{exact-primary-reference}

Shioda, T., \& Katsura, T. (1979). On Fermat varieties. \emph{Tohoku
Mathematical Journal, Second Series, 31}(1), 97--115.
https://doi.org/10.2748/tmj/1178229881 . \textbf{Exact locator:} Remark
1.10, p.~102, equation \((1.1')\), which states that the same
construction applies to smooth \(f=0\) and \(g=0\) in disjoint variable
sets and their three displayed augmented/sum hypersurfaces. The original
rational-map and resolution formulas are Lemmas 1.1--1.2, pp.~98--99.
Primary scan:
https://www.jstage.jst.go.jp/article/tmj1949/31/1/31\_1\_97/\_pdf .

The local primary scan was directly inspected, including printed p.~102.
The full simultaneous induction and the degree/base-ideal verification
above are the deductions checked in this module; the general blowup
cohomology and trace closure inputs are those already checked in
CC-XC08.
